\subsection{2.8 Vergleich verschiedener Projektionsoperatoren}

Mit der Einführung der Projektionstheorie änderte sich der Schwerpunkt der numerischen Untersuchungen grundlegend.

Nicht mehr der OPP selbst stand im Mittelpunkt, sondern die Frage:

\[
\boxed{
\text{Welche Projektion erzeugt eine physikalisch sinnvolle Raumzeit?}
}
\]

Dabei wurde bewusst auf Verfahren zurückgegriffen, die aus der modernen Mathematik und Datenanalyse stammen und keine spezielle physikalische Information enthalten.

Die Hoffnung bestand darin, dass sich eine Projektion als natürlich herausstellt, weil sie unabhängig von den Anfangsbedingungen immer wieder ähnliche Strukturen erzeugt.

%%%%%%%%%%%%%%%%%%%%%%%%%%%%%%%%%%%%%%%%%%%%%%%%%%%%%%%%%%%%%%

\section*{Grundprinzip}

Alle untersuchten Verfahren besitzen dieselbe mathematische Aufgabe:

Aus einer hochdimensionalen relationalen Struktur

\[
R
\]

soll eine niedrigdimensionale Darstellung

\[
Y
\]

erzeugt werden.

Formal:

\[
\mathcal{P}
:
R
\rightarrow
Y.
\]

Die Verfahren unterscheiden sich ausschließlich darin,

welche Eigenschaften der ursprünglichen Struktur möglichst gut erhalten bleiben sollen.

%%%%%%%%%%%%%%%%%%%%%%%%%%%%%%%%%%%%%%%%%%%%%%%%%%%%%%%%%%%%%%

\section*{Die untersuchten Verfahren}

Im Verlauf der numerischen Untersuchungen wurden insbesondere folgende Projektionsoperatoren getestet.

\begin{enumerate}

\item
Principal Component Analysis (PCA)

\item
Kernel Principal Component Analysis (Kernel PCA)

\item
Diffusion Maps

\item
Random Projection

\item
verschiedene Kernel-Familien

\end{enumerate}

Alle Verfahren wurden unter identischen Randbedingungen auf denselben OPP angewendet.

%%%%%%%%%%%%%%%%%%%%%%%%%%%%%%%%%%%%%%%%%%%%%%%%%%%%%%%%%%%%%%

\section*{Principal Component Analysis}

Die PCA stellt den klassischen Standard der Dimensionsreduktion dar.

Sie bestimmt jene Richtungen,

in denen die Varianz maximal ist.

Mathematisch erfolgt dies über die Eigenwertzerlegung einer Kovarianzmatrix.

In den Tests zeigte PCA reproduzierbare Ergebnisse,

jedoch keine vollständige Stabilität.

Typischerweise ergab sich

\[
d_{\mathrm{eff}}
\approx
3.6.
\]

Ein Beispiel aus Test 032:

\[
d_{\mathrm{eff}}
=
3.605
\pm
0.168.
\]

Damit lag PCA bereits überraschend nahe an einer vierdimensionalen Struktur.

%%%%%%%%%%%%%%%%%%%%%%%%%%%%%%%%%%%%%%%%%%%%%%%%%%%%%%%%%%%%%%

\section*{Random Projection}

Als Kontrollversuch wurde eine vollständig zufällige Projektion verwendet.

Sie besitzt keinen Bezug zur Struktur des OPP.

Die Ergebnisse fielen deutlich schlechter aus.

Typischerweise wurde beobachtet

\[
d_{\mathrm{eff}}
\approx
2.9
\]

bei gleichzeitig deutlich größerer Streuung.

Ein Beispiel:

\[
d_{\mathrm{eff}}
=
2.901
\pm
0.415.
\]

Dies zeigte,

dass nicht jede Dimensionsreduktion automatisch sinnvolle Ergebnisse liefert.

%%%%%%%%%%%%%%%%%%%%%%%%%%%%%%%%%%%%%%%%%%%%%%%%%%%%%%%%%%%%%%

\section*{Kernel PCA}

Kernel PCA erweitert die gewöhnliche PCA,

indem zunächst eine nichtlineare Abbildung

\[
\phi
:
x
\rightarrow
\mathcal{H}
\]

in einen hochdimensionalen Hilbertraum erfolgt.

Erst dort wird anschließend eine PCA durchgeführt.

Dadurch können gekrümmte Mannigfaltigkeiten wesentlich besser beschrieben werden.

Die numerischen Resultate waren überraschend stabil.

Ein typisches Ergebnis lautete

\[
d_{\mathrm{eff}}
=
3.910
\pm
0.040.
\]

Dies war der erste Projektionsoperator,

der reproduzierbar Werte praktisch bei vier lieferte.

%%%%%%%%%%%%%%%%%%%%%%%%%%%%%%%%%%%%%%%%%%%%%%%%%%%%%%%%%%%%%%

\section*{Diffusion Maps}

Noch bemerkenswerter entwickelte sich Diffusion Maps.

Dieses Verfahren basiert nicht auf maximaler Varianz,

sondern auf Diffusionsprozessen auf einem Graphen.

Der Kern lautet

\[
P
=
D^{-1}R,
\]

wobei

\[
D
\]

die Gradmatrix bezeichnet.

Die Eigenfunktionen dieses Operators beschreiben großskalige Diffusionsmoden.

Bereits der erste Vergleich ergab

\[
d_{\mathrm{eff}}
=
4.000
\pm
0.000.
\]

Die perfekte Stabilität war zunächst äußerst überraschend.

%%%%%%%%%%%%%%%%%%%%%%%%%%%%%%%%%%%%%%%%%%%%%%%%%%%%%%%%%%%%%%

\section*{Warum dieses Ergebnis kritisch geprüft werden musste}

Gerade wegen der auffallenden Stabilität wurde das Resultat nicht sofort akzeptiert.

Stattdessen wurde untersucht,

ob der Algorithmus selbst bereits eine Vierdimensionalität bevorzugt.

Die Analyse zeigte:

Die verwendeten Diffusionskoordinaten wurden über

\[
K=4
\]

Eigenvektoren konstruiert.

Damit konnte die perfekte Vier

zumindest teilweise algorithmisch bedingt sein.

Das Ergebnis war daher interessant,

aber noch kein Beweis.

%%%%%%%%%%%%%%%%%%%%%%%%%%%%%%%%%%%%%%%%%%%%%%%%%%%%%%%%%%%%%%

\section*{Vergleich der Verfahren}

Die ersten direkten Vergleiche ergaben folgendes qualitative Bild:

\[
\begin{array}{l|c}
\text{Projektionsverfahren}
&
d_{\mathrm{eff}}
\\
\hline
\text{Random Projection}
&
\approx2.9
\\
\text{PCA}
&
\approx3.6
\\
\text{Kernel PCA}
&
\approx3.9
\\
\text{Diffusion Maps}
&
\approx4.0
\end{array}
\]

Bereits dieses Ranking war ausgesprochen aufschlussreich.

%%%%%%%%%%%%%%%%%%%%%%%%%%%%%%%%%%%%%%%%%%%%%%%%%%%%%%%%%%%%%%

\section*{Die erste Auswahl}

Auf Grundlage dieser Ergebnisse wurden zwei Verfahren für weitere Untersuchungen ausgewählt:

\[
\boxed{
\text{Diffusion Maps}
}
\]

und

\[
\boxed{
\text{Kernel PCA}.
}
\]

Alle übrigen Verfahren wurden zunächst zurückgestellt.

Der Grund war einfach.

Beide Methoden lieferten

\begin{itemize}

\item

stabile Resultate,

\item

geringe Standardabweichungen,

\item

Werte nahe der gesuchten Vier.

\end{itemize}

%%%%%%%%%%%%%%%%%%%%%%%%%%%%%%%%%%%%%%%%%%%%%%%%%%%%%%%%%%%%%%

\section*{Physikalische Interpretation}

Erstmals entstand hier eine bemerkenswerte Vermutung.

Vielleicht existiert überhaupt keine eindeutige Projektion.

Stattdessen könnte gelten:

\[
\boxed{
\text{Physikalisch sinnvolle Projektionen erhalten globale Strukturen.}
}
\]

Genau dies leisten sowohl Kernel PCA

als auch Diffusion Maps.

Beide Verfahren versuchen,

nicht lokale Details,

sondern die globale Geometrie des relationalen Netzes zu bewahren.

Dies könnte erklären,

warum gerade sie stabile vierdimensionale Strukturen erzeugen.

%%%%%%%%%%%%%%%%%%%%%%%%%%%%%%%%%%%%%%%%%%%%%%%%%%%%%%%%%%%%%%

\section*{Ein neuer Forschungsansatz}

Nach Abschluss der Vergleichstests wurde die Strategie erneut angepasst.

Anstatt weitere Projektionsverfahren zu testen,

konzentrierte sich die Untersuchung nun auf die beiden stärksten Kandidaten.

Die zentrale Frage lautete jetzt:

\[
\boxed{
\text{Sind die beobachteten Vierdimensionen robust gegenüber Parameteränderungen?}
}
\]

Dies führte unmittelbar zu den umfangreichen Stabilitätsanalysen der Diffusionsprojektion,

welche den Schwerpunkt der folgenden Kapitel bilden.

%%%%%%%%%%%%%%%%%%%%%%%%%%%%%%%%%%%%%%%%%%%%%%%%%%%%%%%%%%%%%%

\section*{Zusammenfassung}

Die systematische Untersuchung verschiedener Projektionsoperatoren führte erstmals zu einer klaren Priorisierung.

Die wichtigsten Erkenntnisse lauten:

\begin{itemize}

\item
Nicht jede Projektion erzeugt physikalisch sinnvolle Strukturen.

\item
Random Projection scheidet als universeller Projektionsoperator weitgehend aus.

\item
PCA liefert bereits erstaunlich gute Ergebnisse.

\item
Kernel PCA erzeugt reproduzierbar Werte nahe
\(
d_{\mathrm{eff}}\approx4.
\)

\item
Diffusion Maps liefert die stabilsten Resultate und wird deshalb zum Hauptkandidaten der weiteren Untersuchungen.

\end{itemize}

Damit begann die zweite Phase der numerischen PST-Forschung:

Nicht mehr der Vergleich verschiedener Verfahren,

sondern die detaillierte Analyse der Diffusionsprojektion selbst.

